% !PRINTED 650
\weberpage{650}

Dann können wir den Satz beweisen:

1. \gesperrt{Ist die Gruppe eines Abel'schen Körpers zerlegbar,
so ist auch der Körper zusammengesetzt.}

Nehmen wir, um dies zu beweisen, an, die Gruppe $\mathfrak{G}$ des
Körpers $\Omega$ sei nach der Formel (1) zerlegbar und $m, a, b$
seien die Grade von $\mathfrak{G}, \mathfrak{A}, \mathfrak{B}$. Wir nehmen eine zur Gruppe $\mathfrak{B}$
\gesperrt{gehörige} Zahl $\xi$ des Körpers $\Omega$, die also durch die Substitutionen
 von $\mathfrak{G}$ in $a$ verschiedene conjugirte Werthe übergeht und
ebenso eine zu $\mathfrak{A}$ gehörige $b$-werthige Zahl $\eta$. Wir können
dann eine Zahl
\[
x = \alpha \xi + \beta \eta
\]
mit \gesperrt{rationalen} Zahlencoëfficienten $\alpha, \beta$ bilden, die $m$ verschiedene
 conjugirte Werthe hat, und dann ist $\Omega = R(x)$. Also
ist $\Omega$ aus $R(\xi)$ und $R(\eta)$ zusammengesetzt (Bd. I, \S.\,143).

Erinnern wir uns nun an die in \S.\,9, 10 dieses Bandes abgeleitete
 Darstellung einer Abel'schen Gruppe durch eine Basis, so
ergiebt sich durch wiederholte Anwendung des eben Bewiesenen:

2. \gesperrt{Jeder Abel'sche Körper $\Omega$ lässt sich aus solchen
zusammensetzen, deren Grad eine Primzahlpotenz
 und deren Gruppe cyklisch ist. Die Grade
dieser zusammensetzenden Körper sind die Invarianten
 der Gruppe von $\Omega$.}

Diese Sätze werden durch den folgenden ergänzt:

3. \gesperrt{Ein Abel'scher Körper von Primzahlpotenzgrad
mit cyklischer Gruppe ist einfach.}

Wenn die Gruppe $\mathfrak{G}$ cyklisch ist, so lassen sich ihre Elemente
 durch die Potenzen einer Basis, $G^\gamma$, darstellen, wenn $\gamma$
ein volles Restsystem nach dem Modul $m$ durchläuft. Wir erhalten
 alle Theiler $\mathfrak{A}$ von $\mathfrak{G}$ dargestellt durch
\[
\mathfrak{A} = G^{b\gamma},
\]
wenn $b$ ein Theiler von $m$ ist und $\gamma$ ein volles Restsystem nach
dem Modul $a = m : b$ durchläuft.

Ist
\[
\mathfrak{A}' = G^{b'\gamma}
\]
ein zweiter Theiler von $\mathfrak{G}$ vom Grade $a'$ und ist
\[
b' \leqq b, \quad a' \geqq a,
\]
so ist, wenn wir jetzt annehmen, dass $m$ und mithin auch
$a, b, a', b'$ Potenzen einer Primzahl $q$ sind, $b'$ ein Theiler von $b$
und folglich $\mathfrak{A}$ ein Theiler von $\mathfrak{A}'$.
